\subsection{Separated roots on the negative quadratic-drift scale}
\label{sec:direct-negative-separated}

Even for separated roots, the limiting endpoint rates can agree.
We compare derivatives of these rates in the common logarithmic height
coordinate. Retaining their drift weights makes the curvature
comparison strict at every interior height.

Fix \(a_0^-,\underline b,\nu>0\), and consider
\begin{equation}
\label{eq:direct-negative-separated-region}
 a\ge a_0^-,\qquad b\ge \underline b,\qquad b^2-4a\ge\nu
\end{equation}
in \eqref{eq:direct-negative-secondary-scales}.
The first positive secondary root is simple:
\[
 \alpha=\frac{b-\sqrt{b^2-4a}}2+O(r),\qquad
 c<\alpha<C,\qquad K_+'(\alpha)<-c.
\]
Set \(x_L=\alpha/4,\ w_a=\alpha/2,\ R_p=3\alpha/4\).
The exact drift and its quadratic comparison have fixed
positive margins up to \(R_p\).

\begin{proposition}[Separated secondary roots]
\label{prop:direct-negative-separated}
In \eqref{eq:direct-negative-separated-region}, \(\Delta_p\)
has at most three zeros for all sufficiently
small independent \(s,r>0\), including parameters at
which the two limiting characteristic rates agree.
\end{proposition}

\begin{proof}
Construct the positive boundary-flow coordinate at a
fixed small output \(z=t_0Z_f,\ t_0=s^2r^2\).
The root-approach interval of central heights is
\begin{equation}
\label{eq:direct-negative-separated-root-cell}
 w=\mathcal W_+(x(t_0Z_f;h),Z_f)\in[w_a,\alpha).
\end{equation}
Advancing from a projected entry at most \(x_L\) to \(w_a\)
requires a fixed positive \(F_+\)-flow time \(\tau_*\).
Thus each central height in this interval satisfies
\[
 h/(s^3r^2)\le (Z_f/s)\e^{-\tau_*/s}.
\]
The regular prefix gives \eqref{eq:direct-common-phase} at the
fixed output \(z=t_0Z_f\). The positive drift margin at \(x_L\)
and the advance \(\tau_*\) verify
\cref{lem:direct-weighted-phase} with \(\delta=s\), \(N=1\) and
\(j=2\). Thus, in \(T=-s\log h\), the phase errors are
\(o(1)F_+(w)\) through two \(T\)-derivatives.
\Cref{lem:boundary-flow-endpoint}, with \(\ell=sr\) and
\(\delta=s\), gives \(-I_{+,w}=E_fA_pw(1+O_{C^3}(s))\), with
\(E_f=t_0(z_*/(t_0Z_f))^{2s}\to0\).
Apply \eqref{eq:direct-rate-curvature}:
\[
 \mathfrak H_+'=
 F_+(w)\{Q_+(w)+o(1)\},\qquad
 Q_+(w)=\frac{wK_+''-K_+'}{w^2}.
\]
By \eqref{eq:direct-curvature-coefficient-cubic}, with \(K_+'(0)=-b\) and
\(R_K(x)=-rR_q(srx)\), one has \(Q_+(w)=b/w^2+O(r)>c_Q>0\), since
\(b\ge \underline b\) and \(w<\alpha<C\).
Consequently \(\mathfrak H_+'\ge cF_+(w)>0\) on the root interval.

The negative output at the same central height may lie on either
scale. We follow it from bounded secondary outputs through the annulus
and the compact primary interval to its far root, retaining the
endpoint estimates needed to absorb the exterior curvature in each
range.
Choose a small fixed \(x_c>0\) so that the negative boundary action
to \(2x_c\) is less than \(\tau_*/4\). This is possible
because the central coefficient is at least \(a_0^-\).
The negative prefix coordinate gives exponentially small height
at this cut.
The negative secondary drift is
\[
 K_b(X)=a+bX+X^2+rX^3R_q(-srX),\qquad F_b=K_b/X .
\]
Choose a large fixed \(M\), and a small fixed output
height \(Z_b>0\). The boundary-flow coordinate on
a box containing \([x_c/4,5M]\) shows that every
negative output with \(X\le4M\) at \(t_0Z_b\)
has projected state \(w_b\in[x_c,4M]\).
At the regular entry \(x_c/2\), the drift \(F_b\) has a fixed
positive lower bound. The advance from that entry to
\(w_b\ge x_c\), at the fixed output \(z=t_0Z_b\), has a fixed
positive phase gap. Hence \cref{lem:direct-weighted-phase}
applies with \(\delta=s\), \(N=1\) and \(j=2\), in the same
\(T\), giving phase errors \(o(1)F_b(w_b)\) through two
\(T\)-derivatives.
By \cref{lem:boundary-flow-endpoint}, its endpoint has amplitude
\(E_b=t_0(z_*/(t_0Z_b))^{2s}\to0\). By
\eqref{eq:direct-curvature-coefficient-cubic}, with \(K_b'(0)=b\) and
\(R_K(X)=rR_q(-srX)\), one has \(Q_b(X)=-b/X^2+O(r)<-c_M<0\) for
\(x_c\le X\le4M\). The endpoint rate identities and the weighted phase
estimate then give
\[
 \mathfrak H_-'\le-cF_b(w_b),\qquad
 |I_{-,T}|+|I_{-,TT}|\le CE_bF_b(w_b).
\]

If the radius \(X\) of the negative graph at \(t_0Z_b\) exceeds
\(4M\), the graph is thin at \(X=M\). Indeed, the outward graph then
reaches \(X=4M\) below \(Z_b\), it lies in the coordinate box between
\(X=M\) and \(X=4M\), and \(8B_0Z_b\le15M^2\) since \(Z_b\) is small.
For small \(s\), \cref{lem:direct-radial-advance} with \(\delta=s\),
\(X_1=M\), \(X_2=4M\) and \(C_F=\sup|F_b|\) on the box bounds its
secondary height at \(X=M\) by
\(Z_b\e^{-3M/(2C_Fs)}<sM^2\delta_{\rm an}\).

For a thin exit in \(Mr\le\rho\le v_0\), the negative regular
prefix allows us to apply \cref{lem:direct-annulus-rates}, with
\(u=r/\rho\), \(v=\rho\). Choose \(M\) large and \(v_0,s\) small.
Since \(b\ge \underline b\), \eqref{eq:direct-annulus-clock-rate} gives
\[
 \mathfrak H_-'\le-c(u+v),\qquad
 |I_{-,T}|+|I_{-,TT}|\le C(s\rho)^{2-4s}.
\]
For \(v_0/2\le\rho\le0.9\),
\cref{lem:negative-branch-far-root}(i) instead gives
\(\mathfrak H_-'<-c_{v_0}<0\), with the same
endpoint size estimate.

Every graph still thin at \(0.9\) satisfies
\eqref{eq:direct-negative-far-core}, and
\cref{lem:negative-branch-far-root}(ii) gives, uniformly up to the
far root,
\[
 \mathfrak H_-'\le-cF_n(w_n),\qquad
 |I_{-,T}|+|I_{-,TT}|\le CE_nF_n(w_n).
\]
These ranges cover every negative-branch output. We now compare the
exterior term with the curvature margin in each range. Equation
\eqref{eq:direct-exterior-curvature} gives exterior contributions
\(o(1)F_b\) and \(o(1)F_n\) on the bounded secondary and far-root
intervals and \(o(1)\) on the compact primary interval. On the small
primary annulus, \eqref{eq:direct-annular-exterior-curvature} gives
\(o(\rho)\).
Each exterior contribution is absorbed by its negative curvature
margin. Thus
\[
 \mathcal L''=\mathfrak H_-'-\mathfrak H_+'
                      +\partial_T^2[(\log\Pi_p')(I_-)]<0
\]
throughout \eqref{eq:direct-negative-separated-root-cell}.
The positive drift weights give strict negativity at every interior
central height, although both may tend to zero at the endpoint.
Thus the curvature remains strictly negative at each interior height
even when the limiting rates agree. In the leading model the
characteristic rates agree at
\(\alpha=b/3,\ a=2b^2/9\), which is included here.

Finally apply \cref{lem:prefix-root-cover} with \(\delta=s\),
\(\ell=sr\), \(\epsilon_3=r\), comparison drifts \(a-\sigma bx+x^2\)
and the radii \(w_a,R_p,\alpha\) fixed above. The positive margin
and the gap \(R_p-w_a\) are fixed, and the inverse boundary-flow
coordinate is defined up to \(\alpha\) at \(Z_f\).
The root range is \eqref{eq:direct-negative-separated-root-cell},
where \(\mathcal L_{TT}<0\). Hence \(\Delta_p\) has at most three
zeros on \(\cH_p\).
\end{proof}

\subsubsection{The compact secondary coefficient cover}

For any fixed \(a_0^->0\), the normalized coefficients
with \(a\ge a_0^-\) now have a bound of three.
Choose the fixed fold width \(\epsilon_{\rm f}\) from
\cref{prop:direct-negative-fold}, and let \(\nu=\epsilon_{\rm f}^2/50\).
The three overlapping regions are
\[
 b^2-4a\le-\nu,\qquad |b^2-4a|\le2\nu,\qquad b^2-4a\ge\nu .
\]
They are covered respectively by
\cref{prop:direct-buffered-negative,prop:direct-negative-fold,prop:direct-negative-separated}.
For the last two regions normalization forces \(b=1\):
otherwise \(a=1\) would give \(b^2-4a\le-3\).
On the middle region
\mbox{\(|a-b^2/4|\le\epsilon_{\rm f}^2/100\)}, as required for the fold.
Each of the three propositions bounds the zeros on the whole height
interval. The minimum of their scale thresholds therefore gives a
uniform bound of three on this coefficient set.
